% GENERATED FILE. Edit canonical lessons, then run npm run generate:books.
% canonical-source-hash: fnv1a64:db2c6eb6817efe8b
% canonical-lessons: LA-C517-vena, LA-C517-nervus, LA-C517-pellis, LA-C517-sudor, LA-C517-pulsus

\chapter{Inside the body}
\label{ch:la-inside-the-body}
\hlchaptermodality{517}

\begin{chapteropening}
\textbf{By the end of this chapter:} \emph{I can talk about veins, sinews, skin and sweat.}

Complete the last lesson of chapter 517: I can talk about veins, sinews, skin and sweat.
\end{chapteropening}

\section[vēna, vēnae]{vēna, vēnae — a vein}

\label{lesson:LA-C517-vena}



\begin{quote}
Before the new one: say the Latin for a limb, then the Latin for joints.
\end{quote}

\subsection*{You'll want to know: vēna, vēnae}
\textbf{vēna, vēnae} — \textquotedblleft{}a vein\textquotedblright{}.

\begin{quote}
Sanguis per vēnās currit.
\end{quote}

\begin{quote}
\emph{Blood runs through the veins.}
\end{quote}

Feminine, first declension. Also a vein of metal, or of talent. English \emph{vein}.

\begin{cousinweb}
\textbf{vēna} is a native Latin noun; Lewis \& Short guesses a link with \emph{vehō} \textquotedblleft{}to carry\textquotedblright{}, but the origin is not settled. English \emph{vein} comes from it through French.
\end{cousinweb}

\begin{grammarlens}[title={What you've built}]
One more word for telling, explaining and describing at B1.
\end{grammarlens}

\subsection*{Your turn}
\begin{itemize}
\raggedright
  \item \emph{Say it:} \emph{vēna, vēnae}
  \item \emph{Say it:} \emph{vēna, vēnae}, once more
  \item \emph{Say it:} a sentence of your own with \emph{vēna, vēnae}
  \item \emph{Recall:} say the Latin for a limb, then the Latin for joints, then say \emph{vēna, vēnae} again
\end{itemize}

\begin{tcolorbox}[breakable,colback=teal!4,colframe=teal!35!black,title={Before you move on}]
What does vēna, vēnae mean? (\textquotedblleft{}A vein\textquotedblright{}.) Say it once more.
\end{tcolorbox}
\section[nervus, nervī]{nervus, nervī — a sinew, tendon; a string}

\label{lesson:LA-C517-nervus}



\begin{quote}
Before the new one: say the Latin for joints, then the Latin for a vein.
\end{quote}

\subsection*{You'll want to know: nervus, nervī}
\textbf{nervus, nervī} — \textquotedblleft{}a sinew, tendon; a string\textquotedblright{}.

\begin{quote}
Faber nervōs firmōs habet.
\end{quote}

\begin{quote}
\emph{The smith has strong sinews.}
\end{quote}

Masculine, second declension. In the plural also strength, vigour. English \emph{nerve}.

\begin{cousinweb}
\textbf{nervus} is inherited from a Proto-Indo-European word for \textquotedblleft{}sinew, tendon\textquotedblright{}, and Greek \emph{neuron} is a cousin (EtymDB; Lewis \& Short). English \emph{nerve} comes from it.
\end{cousinweb}

\begin{grammarlens}[title={What you've built}]
One more word for telling, explaining and describing at B1.
\end{grammarlens}

\subsection*{Your turn}
\begin{itemize}
\raggedright
  \item \emph{Say it:} \emph{nervus, nervī}
  \item \emph{Say it:} \emph{nervus, nervī}, once more
  \item \emph{Say it:} a sentence of your own with \emph{nervus, nervī}
  \item \emph{Recall:} say the Latin for joints, then the Latin for a vein, then say \emph{nervus, nervī} again
\end{itemize}

\begin{tcolorbox}[breakable,colback=teal!4,colframe=teal!35!black,title={Before you move on}]
What does nervus, nervī mean? (\textquotedblleft{}A sinew, tendon; a string\textquotedblright{}.) Say it once more.
\end{tcolorbox}
\section[pellis, pellis]{pellis, pellis — a hide, skin}

\label{lesson:LA-C517-pellis}



\begin{quote}
Before the new one: say the Latin for a vein, then the Latin for a sinew.
\end{quote}

\subsection*{You'll want to know: pellis, pellis}
\textbf{pellis, pellis} — \textquotedblleft{}a hide, skin\textquotedblright{}.

\begin{quote}
Pāstor pellem ovis vendit.
\end{quote}

\begin{quote}
\emph{The shepherd sells the sheep's hide.}
\end{quote}

Feminine, third declension. Mostly an animal's hide; \emph{cutis} is skin as part of the body.

\begin{cousinweb}
\textbf{pellis} is inherited from a Proto-Indo-European root meaning \textquotedblleft{}to cover, skin\textquotedblright{} (EtymDB), and Old English \emph{fell} \textquotedblleft{}hide\textquotedblright{} is a cousin. English \emph{pellicle}, a thin skin, comes from its diminutive.
\end{cousinweb}

\begin{grammarlens}[title={What you've built}]
One more word for telling, explaining and describing at B1.
\end{grammarlens}

\subsection*{Your turn}
\begin{itemize}
\raggedright
  \item \emph{Say it:} \emph{pellis, pellis}
  \item \emph{Say it:} \emph{pellis, pellis}, once more
  \item \emph{Say it:} a sentence of your own with \emph{pellis, pellis}
  \item \emph{Recall:} say the Latin for a vein, then the Latin for a sinew, then say \emph{pellis, pellis} again
\end{itemize}

\begin{tcolorbox}[breakable,colback=teal!4,colframe=teal!35!black,title={Before you move on}]
What does pellis, pellis mean? (\textquotedblleft{}A hide, skin\textquotedblright{}.) Say it once more.
\end{tcolorbox}
\section[sūdor, sūdōris]{sūdor, sūdōris — sweat}

\label{lesson:LA-C517-sudor}



\begin{quote}
Before the new one: say the Latin for a sinew, then the Latin for a hide.
\end{quote}

\subsection*{You'll want to know: sūdor, sūdōris}
\textbf{sūdor, sūdōris} — \textquotedblleft{}sweat\textquotedblright{}.

\begin{quote}
Sūdor dē fronte agricolae fluit.
\end{quote}

\begin{quote}
\emph{Sweat flows from the farmer's brow.}
\end{quote}

Masculine, third declension, from \emph{sūdō} — to sweat.

\begin{cousinweb}
\textbf{sūdor} is built on \emph{sūdō} \textquotedblleft{}to sweat\textquotedblright{}, inherited from a Proto-Indo-European root meaning \textquotedblleft{}to sweat\textquotedblright{}; English \emph{sweat} is a cousin (EtymDB).
\end{cousinweb}

\begin{grammarlens}[title={What you've built}]
One more word for telling, explaining and describing at B1.
\end{grammarlens}

\subsection*{Your turn}
\begin{itemize}
\raggedright
  \item \emph{Say it:} \emph{sūdor, sūdōris}
  \item \emph{Say it:} \emph{sūdor, sūdōris}, once more
  \item \emph{Say it:} a sentence of your own with \emph{sūdor, sūdōris}
  \item \emph{Recall:} say the Latin for a sinew, then the Latin for a hide, then say \emph{sūdor, sūdōris} again
\end{itemize}

\begin{tcolorbox}[breakable,colback=teal!4,colframe=teal!35!black,title={Before you move on}]
What does sūdor, sūdōris mean? (\textquotedblleft{}Sweat\textquotedblright{}.) Say it once more.
\end{tcolorbox}
\section[pulsus, pulsūs]{pulsus, pulsūs — a beat, pulse}

\label{lesson:LA-C517-pulsus}



\begin{quote}
Before the new one: say the Latin for a hide, then the Latin for sweat.
\end{quote}

\subsection*{You'll want to know: pulsus, pulsūs}
\textbf{pulsus, pulsūs} — \textquotedblleft{}a beat, pulse\textquotedblright{}.

\begin{quote}
Medicus pulsum vēnārum temptat.
\end{quote}

\begin{quote}
\emph{The doctor tests the beat of the veins.}
\end{quote}

Masculine, fourth declension, from \emph{pellō} — to strike, push. \emph{Pulsus vēnārum} — the pulse. English \emph{pulse}.

\begin{cousinweb}
\textbf{pulsus} is the action noun of \emph{pellō} \textquotedblleft{}to push, strike\textquotedblright{} (Lewis \& Short; EtymDB). English \emph{pulse} and \emph{repulse} come from the family.
\end{cousinweb}

\begin{grammarlens}[title={What you've built}]
One more word for telling, explaining and describing at B1.
\end{grammarlens}

\subsection*{Your turn}
\begin{itemize}
\raggedright
  \item \emph{Say it:} \emph{pulsus, pulsūs}
  \item \emph{Say it:} \emph{pulsus, pulsūs}, once more
  \item \emph{Say it:} a sentence of your own with \emph{pulsus, pulsūs}
  \item \emph{Recall:} say the Latin for a hide, then the Latin for sweat, then say \emph{pulsus, pulsūs} again
\end{itemize}

\begin{tcolorbox}[breakable,colback=teal!4,colframe=teal!35!black,title={Before you move on}]
What does pulsus, pulsūs mean? (\textquotedblleft{}A beat, pulse\textquotedblright{}.) Say it once more.
\end{tcolorbox}
